% 1. Preamble
\documentclass[12pt, a4paper]{article}

% Import Packages
\usepackage[utf8]{inputenc}        % Encoding
\usepackage{amsmath, amssymb}     % Advanced math formulas
\usepackage{mathtools}           % Math enhancements
\usepackage{graphicx}            % Image insertion
\usepackage{hyperref}            % Interactive hyperlinks & PDF bookmarks
\usepackage{color}               % Color text and backgrounds
\usepackage{geometry}            % Page geometry
\usepackage{comment}             % Multi-line comments
\usepackage{titling}             % Custom title formatting

\geometry{a4paper, margin=1in}

% Custom Command Definitions
\newcommand{\Sball}{\mathcal{S}}
\providecommand{\norm}[1]{\left\lVert#1\right\rVert}

% Hyperref Setup for Clean PDF Links
\hypersetup{
    colorlinks=true,
    linkcolor=blue,
    urlcolor=blue,
    pdftitle={Pattern Arithmetic},
    pdfauthor={Sandeep Lakshminarayan Chiluveru}
}

% Metadata & Custom Title Banner
\pretitle{\begin{center}\Huge\fontfamily{tt}\selectfont ================================================================================\\[0.5em]}
\posttitle{\\[0.5em] ================================================================================\end{center}}

\title{Pattern Arithmetic: Phase-Colored Geometric Framework for High-Density Information via Unit Probability Spheres}
\author{Sandeep Lakshminarayan Chiluveru}
\date{\today}

\begin{document}

\maketitle

% ABSTRACT
\begin{abstract}
Traditional matrix representations of continuous state spaces suffer from severe structural limitations: rigid coordinate axis dependencies, exponential spatial scaling costs ($O(N^3)$), and the destructive overwrite of overlapping signals at identical spatial coordinates. To overcome these constraints, this paper introduces \textit{Pattern Arithmetic}, a continuous geometric framework that models multi-variable information states using state vectors on a Unit Probability Sphere $\mathcal{S} \subset \mathbb{R}^3$, augmented by an internal phase-color degree of freedom $\chi \in [0, 2\pi)$. By mapping continuous probability distributions over vector fields and treating color as an independent phase factor, overlapping states occupy identical spatial coordinates without loss of information density. We establish six foundational operational axioms, including Information Idempotence ($1 \oplus 1 = 1$), Incremental Subsumption ($A + A_{0.01} = A_{0.01}$), and Phase Interference ($A - A = 0$). Finally, we demonstrate how this pattern-based algebra provides a parameter-efficient alternative for high-density latent storage, signal unmixing, and physical field modeling.
\end{abstract}

\section{Introduction}

Ordinary arithmetic uses the integer $1$ as a unit of truth.
The identity $1=1$ says that two tokens are of the same kind.
The rule $1+1=2$ says that combining two such tokens produces a
larger quantitative state. That algebra answers the question
\emph{how many?}

It is not the only algebra. A color and a musical note are also
units of truth, but their combination does not enlarge a pile.
Under additive light, red and green yield yellow. Two copies of
the same yellow, once brightness is normalized, remain yellow.
Two notes yield a chord, not a larger integer; a tone plus its
opposite phase yields silence. In each case the word ``add''
names a different rule, and therefore a different kind of unit.

When the object of interest is a structured whole rather than a
count, a hue, or a pitch, the natural unit is a \emph{pattern}.
This paper represents one pattern by one state on a unit
probability sphere $\mathcal{S} \subset \mathbb{R}^3$:
\begin{itemize}
    \item $r \in [0, 1]$ is the strength of the unit,
          the descendant of the numeral $1$, capped so that
          a state remains a single unit of truth;
    \item $(\theta, \phi)$ is the direction of the pattern
          on the sphere;
    \item $\chi \in [0, 2\pi)$ is an internal phase-color,
          so two identities may occupy the same direction
          without overwriting one another.
\end{itemize}
Pattern Arithmetic is the algebra of these units. Identical
states do not double after normalization ($A \oplus A = A$).
Opposite phases cancel ($A - A = 0$). Rotations of the sphere
preserve the unit.

A single numerical picture is enough to see the difference
from counting. Take two unit states at the same place,
one with phase $0^\circ$ (red) and one with phase $120^\circ$
(green). Counting would return $2$. Phase-aware combination
returns a single unit whose blended phase is $60^\circ$
(yellow). The same two strengths with opposite phase
($0^\circ$ and $180^\circ$) return the null state, black.
Two units of truth become either one new identity or none.

\subsection{The Limitations of Discretized Matrix Grids}
A discrete matrix grid is the wrong box for such a unit. Traditional matrix representations of continuous state spaces suffer from two severe structural bottlenecks:
\begin{itemize}
    \item \textbf{Memory Scaling Bottlenecks:} A voxel cube of resolution $N$ stores $O(N^3)$ cells, growing steeply with resolution.
    \item \textbf{Signal Identity Destruction:} Mapping multiple overlapping variables to identical spatial coordinates often results in the destructive overwriting of underlying signal identities.
\end{itemize}

\subsection{The Phase-Colored Unit Probability Sphere Framework}
To resolve these grid constraints, the sphere is offered as one bounded object that can carry overlapping identities by direction and phase, without enlarging a lattice.
\begin{itemize}
    \item \textbf{Spatial Orientation $(\theta, \phi)$:} Defines continuous directional coordinates across the spherical surface.
    \item \textbf{Magnitude Likelihood ($0 \leq r \leq 1$):} Represents bounded probability density within the sphere.
    \item \textbf{Internal Color Phase ($\chi \in [0, 2\pi)$):} Serves as an independent phase degree of freedom.
\end{itemize}

\subsection{Core Contributions \& High-Level Intuition}
The principal contributions of Pattern Arithmetic are:
\begin{itemize}
    \item \textbf{High-Density Superposition:} Storing overlapping state vectors concurrently within a single bounded spatial manifold.
    \item \textbf{Parameter-Efficient Field Storage:} Eliminating discrete voxel bounds by treating states as continuous probability distributions.
\end{itemize}

The next section turns these definitions into formal mathematical axioms.


\section{Axiomatic Foundations Of Pattern Arithmetic}

\subsection{Axiom 1: The Normalized Unit Truth State ($\mathcal{S} = 1$)}
\label{subsec:axiom1}
To ensure stability across multi-variable operations, every state vector field $\vec{v} = (r, \theta, \phi, \chi)$ on the Unit Probability Sphere $\mathcal{S} \subset \mathbb{R}^3$ represents an immutable unit of truth ($\mathcal{S} = 1$). Global likelihood integrates strictly to unity:
\begin{equation}
    \iiint_{\mathcal{S}} P(r, \theta, \phi) \, dV = 1.0
    \label{eq:unit_truth}
\end{equation}

\subsection{Axiom 2: The Absolute Null State ($\mathcal{S} = 0$, Black)}
\label{subsec:axiom_null}
The state $\mathcal{S} = 0$ defines the absolute lower bound of the state space, corresponding to complete destructive interference, total phase cancellation, or zero-state equilibrium:
\begin{equation}
    \norm{\vec{v}} = 0 \quad \forall \; (\theta, \phi, \chi) \implies \mathcal{S} = 0 \quad \text{(Black)}
    \label{eq:null_state}
\end{equation}
Physically and visually, $\mathcal{S} = 0$ represents an empty probability field devoid of active signal energy or directional orientation.

\subsection{Axiom 3: The Infinite Superposition State ($\mathcal{S} = \infty$, White)}
\label{subsec:axiom_infinite}
The state $\mathcal{S} = \infty$ defines the theoretical upper bound of maximal entropy, representing total spectral superposition across all continuous orientations $(\theta, \phi)$ and internal color-phase frequencies $\chi \in [0, 2\pi)$:
\begin{equation}
    \mathcal{S} = \lim_{N \to \infty} \bigoplus_{i=1}^{N} A_i \implies \mathcal{S} = \infty \quad \text{(White)}
    \label{eq:infinite_state}
\end{equation}
Physically and visually, $\mathcal{S} = \infty$ corresponds to full-spectrum white light, where all possible state vectors co-exist simultaneously without spatial exclusion.

\subsection{Axiom 4: Information Idempotence ($1 \oplus 1 = 1$)}
\label{subsec:axiom_idempotence}
Combining identical information states yields no expansion of likelihood volume. Upon intensity normalization:
\begin{equation}
    A \oplus A = \frac{A + A}{\norm{A + A}} = A
    \label{eq:idempotence_formal}
\end{equation}

\subsection{Axiom 5: Phase-Aware Superposition \& Interference}
\label{subsec:axiom_interference}
The combination of two truth states $A$ and $B$ is governed by relative phase difference $\Delta\chi$:
\begin{equation}
    A \pm B = \sqrt{\norm{A}^2 + \norm{B}^2 + 2\norm{A}\norm{B}\cos(\Delta\chi)}
    \label{eq:interference_formal}
\end{equation}
Destructive anti-phase interference ($\Delta\chi = 180^\circ$) collapses composite states directly into the Null State ($A - A = 0$, Black) defined in Axiom~\ref{subsec:axiom_null}.

\subsection{Axiom 6: Rotational Conservation (\texorpdfstring{$\times \mathbf{R}$}{x R})}
\label{subsec:axiom_rotation}
Multiplication operates as a Lie group rotation $\mathbf{R} \in \mathrm{SO}(3)$ in continuous spherical space, preserving total truth volume:
\begin{equation}
    \norm{\mathbf{R} \cdot A} = \norm{A} = 1.0
\end{equation}

\section{Transformation Algebra: From Identity to Expression}
\label{sec:transformation_algebra}

\subsection{Motivation: Two Layers of Meaning}
\label{subsec:transform_motivation}
The operators established so far ($\oplus$, $\pm$) answer the question
\emph{what is this state?} Combining Sky, Sun, and Wind under $\oplus$
yields a single, order-independent resultant --- a new composite unit
of truth, exactly as combining Ruler and Man yields the composite unit
King. This is a \emph{lexical} operation: it mints identities, the way
a dictionary mints words from morphemes. It does not, and structurally
cannot, distinguish a poem in which the sun rises through a wind-blown
sky from a poem in which wind extinguishes a darkening sun --- both
share the identical resultant
$R = \text{Sky} \oplus \text{Sun} \oplus \text{Wind}$, since $\oplus$
is commutative and idempotent by Axioms~\ref{subsec:axiom_idempotence}
and~\ref{subsec:axiom_interference}.

The distinction between these two poems is not a difference in
\emph{ingredients}; it is a difference in \emph{how one state acts
upon another, and in what order}. This is a \emph{relational}
operation, and it requires a second operator, $T$, layered on top of
$\oplus$: where $\oplus$ builds nouns, $T$ builds sentences. A poem is
not the resultant $R$ itself, but a specific trajectory traced across
$\mathcal{S}$ by successive transformations applied to $R$.

\subsection{Formal Definition of the Transformation Operator}
\label{subsec:transform_definition}
Rather than introduce $T$ as an unrelated new primitive, it is defined
using the same rotational machinery already given in
Axiom~\ref{subsec:axiom_rotation}: every state $w \in \mathcal{S}$
acts on another state $x \in \mathcal{S}$ by rotating $x$'s direction
about $w$'s own axis, by an angle proportional to $w$'s strength, while
additively shifting $x$'s phase by $w$'s phase:
\begin{equation}
    T_w(x) \;=\; \Big(r_x,\;\; \mathrm{Rot}_{\hat n(\theta_w,\phi_w)}\!\big(\alpha(r_w)\big)\cdot(\theta_x,\phi_x),\;\; (\chi_x + \chi_w) \bmod 2\pi\Big)
    \label{eq:transform_def}
\end{equation}
where $\hat n(\theta_w,\phi_w)$ is the unit axis defined by $w$'s
direction, $\mathrm{Rot}_{\hat n}(\alpha) \in \mathrm{SO}(3)$ is the
rotation of Axiom~\ref{subsec:axiom_rotation} about that axis by angle
$\alpha$, and $\alpha(r_w) = \pi r_w$ is the simplest monotonic choice
mapping strength to rotation angle (a full-strength unit, $r_w = 1$,
rotates by the maximal $180^\circ$). Magnitude $r_x$ is left invariant:
a transformation reorients and recolors a state, it does not
strengthen or weaken it --- that remains the role of $\oplus$.

\subsection{Derived Properties}
\label{subsec:transform_properties}
The following properties follow from Equation~\eqref{eq:transform_def}
rather than being separately postulated.

\paragraph{Identity.} If $w = \mathbf{0}$ (the Null State of
Axiom~\ref{subsec:axiom_null}), then $\alpha(0) = 0$ and the phase
shift is $0$, so $T_{\mathbf{0}} = \mathrm{id}$.

\paragraph{Invertibility.} Define $w^{*} = (r_w, \theta_w, \phi_w, -\chi_w)$.
Then $T_{w^{*}}\!\big(T_w(x)\big) = x$ for all $x$: the same axis and
angle undo the rotation, and the phase shift cancels. Every
transformation has a well-defined inverse. Note that this is distinct
from anti-phase interference ($A - A = 0$, Axiom~\ref{subsec:axiom_interference}),
which annihilates a state; $T_{w^{*}}$ instead \emph{undoes a
transformation} on a state that persists.

\paragraph{Associativity.} Since each $T_w$ is a function
$\mathcal{S} \to \mathcal{S}$, composition of transformations is
associative automatically, by the associativity of function
composition; no additional axiom is required.

\paragraph{Non-commutativity.} Rotations about distinct axes in
$\mathrm{SO}(3)$ generically do not commute: for $a, b$ with distinct
directions $(\theta,\phi)$,
\begin{equation}
    T_a\big(T_b(x)\big) \;\neq\; T_b\big(T_a(x)\big).
    \label{eq:noncommute}
\end{equation}
This is not an additional assumption; it is inherited directly from
the well-established non-commutativity of $\mathrm{SO}(3)$, already
invoked by Axiom~\ref{subsec:axiom_rotation}. It is this property,
and only this property, that allows two poems built from the same
resultant $R$ to diverge: order of application is meaningful even
though the underlying identity is not.

\subsection{Poems as Transformation Trajectories}
\label{subsec:transform_trajectories}
A poem $\pi$ built from resultant $R$ and an ordered sequence of
transforming states $w_1, w_2, \ldots, w_n$ is the trajectory
\begin{equation}
    \pi = \big(R;\, w_1, \ldots, w_n\big): \quad
    R \;\longrightarrow\; T_{w_1}(R) \;\longrightarrow\; T_{w_2}T_{w_1}(R) \;\longrightarrow\; \cdots \;\longrightarrow\; T_{w_n}\cdots T_{w_1}(R).
    \label{eq:poem_trajectory}
\end{equation}
Two poems may share an identical resultant $R$ (the same theme, the
same ``ingredients'') while tracing entirely different trajectories
--- this is the formal sense in which many poems arise from the same
units of truth. Conversely, because rotation composition is
many-to-one, two distinct trajectories may terminate at the same
endpoint; a poem's identity is therefore carried by its full path
through $\mathcal{S}$, not by its final state alone.

\subsection{Worked Numerical Example: Sun, Wind, and Sky}
\label{subsec:transform_example}
Let Sky be the state acted upon, with direction
$(\theta,\phi) = (45^\circ, 45^\circ)$, i.e.\ Cartesian axis
$(0.500, 0.500, 0.707)$. Let Sun define an axis along
$(\theta,\phi) = (90^\circ, 0^\circ)$, i.e.\ $(1,0,0)$, with strength
$r_{\text{Sun}} = 1.0$, giving rotation angle $\alpha_{\text{Sun}} =
180^\circ$. Let Wind define an axis along $(\theta,\phi) = (90^\circ,
90^\circ)$, i.e.\ $(0,1,0)$, with strength $r_{\text{Wind}} = 0.7$,
giving $\alpha_{\text{Wind}} = 126^\circ$.

\emph{Order 1 --- Sun acts first, then Wind} (the sun rising through a
wind-blown sky):
\begin{equation}
    T_{\text{Wind}}\big(T_{\text{Sun}}(\text{Sky})\big) \;\approx\; (-0.866,\, -0.500,\, 0.011)
\end{equation}

\emph{Order 2 --- Wind acts first, then Sun} (wind extinguishing a
darkening sun):
\begin{equation}
    T_{\text{Sun}}\big(T_{\text{Wind}}(\text{Sky})\big) \;\approx\; (0.278,\, -0.500,\, 0.820)
\end{equation}

The two endpoints are clearly distinct, confirming
Equation~\eqref{eq:noncommute}: identical primitives, identical
resultant $R$, opposite structural readings. Notably, if $\chi_{\text{Sun}}
= 0^\circ$ and $\chi_{\text{Wind}} = 90^\circ$, the accumulated phase
shift is $270^\circ$ in \emph{both} orderings, since phase addition is
commutative. The two poems therefore share the same blended color/mood
even as their spatial --- structural --- meaning diverges. This
asymmetry (shared affect, divergent structure) is a direct and
testable prediction of the model.

\subsection{Scope}
\label{subsec:transform_scope}
$T$ is deliberately left unconstrained beyond
Equation~\eqref{eq:transform_def}: any sequence of any transforming
states is a formally valid trajectory. Pattern Arithmetic in this form
supplies the substrate on which expression is built --- the
combinatorial space of possible trajectories --- and does not yet
supply a grammar constraining which trajectories are coherent,
well-formed, or meaningful to a human reader. Whether such constraints
belong inside this framework (as further axioms on admissible
sequences of $T$) or outside it (as a learned or externally specified
grammar layered on top) is left open.

\section{Operational Algebra \& Numerical Demonstration}

\subsection{Vector Addition \& Phase Blending}
Using Equation~\eqref{eq:interference_formal} from Axiom~\ref{subsec:axiom_interference}, combining Red Vector ($0^\circ$) and Blue Vector ($120^\circ$) at coordinate ($0.6, 0.8$) results in deterministic phase blending ($\chi = 66.6^\circ$).

\subsection{Rotational Transposition (Multiplication) \& Unmixing (Division)}
\begin{itemize}
    \item \textbf{Multiplication ($\times$):} State transformation via Lie group rotation matrix $\mathbf{R} \in \mathrm{SO}(3)$.
    \item \textbf{Division ($\div$):} Deconvolution and signal unmixing via inverse transformation $\mathbf{R}^{-1}$ or anti-phase isolation.
\end{itemize}

\subsection{Computational Efficiency \& Spectral Compression}
\begin{itemize}
    \item Mitigating $O(N^3)$ voxel grid costs using Spherical Harmonics ($Y_l^m$) coefficient arrays $O(L^2)$.
\end{itemize}

\section{Real-World Physical \& Biological Models}
\subsection{Forest Ecosystem Dynamics}
\begin{itemize}
    \item Modeling flora, fauna, soil, and moisture resource completion via vector addition ($A + B = AB$).
\end{itemize}

\subsection{Targeted Viral/Pathogen Cancellation (Selective Subtraction)}
\begin{itemize}
    \item Isolating disease signatures and applying $180^\circ$ anti-phase mirror vectors to achieve zero-state equilibrium ($A - B = 0$).
\end{itemize}

\subsection{Atmospheric Storm Systems (Phase Shifts and Shear Turbulence)}
\begin{itemize}
    \item In-phase wave coherence vs. minor phase shift ($0.01$) driving localized vortices and trajectory shifts.
\end{itemize}

\section{Discussion \& Related Work}
\subsection{Comparison to Quantum Bloch Spheres and Hilbert Spaces}
\subsection{Integration with machine learning latent embeddings (Cosine Similarity)}
\subsection{Hardware Realization: Optical Phase Computing vs. Digital GPUs}

\section{Conclusion \& Future Directions}

\end{document}